\section{Conclusiones}

\begin{enumerate}
\item El problema de cobertura de antenas se formuló como un modelo de programación lineal entera binaria con 500 variables y 500 restricciones de cobertura, que minimiza el costo total de instalación y exige que cada cliente sea cubierto por al menos una antena.
\item Para esta instancia densa de $500\times500$, el solver exacto HiGHS no logró demostrar la optimalidad en 600 segundos: terminó con una solución de costo 50\,119 y una brecha de 38,5\,\% respecto a su mejor cota inferior. Esto ilustra por qué el Set Cover, por ser NP-difícil, es costoso de resolver de forma exacta cuando crece el tamaño de la instancia.
\item El Algoritmo Genético con reparación diseñado encontró soluciones factibles en pocos segundos. Su mejor resultado (48\,878) fue 2,5\,\% más barato que la mejor solución del solver en el tiempo límite, y 6,1\,\% más barato que la heurística voraz. El AG es entonces una alternativa práctica cuando el tiempo de cómputo es limitado.
\item El AG es estocástico: el costo varía de una corrida a otra (desviación de unos 1\,500 en la configuración base). Aumentar la población o usar una búsqueda más larga reduce esa variabilidad (desviación de 606 en la configuración extendida), a cambio de más tiempo.
\item El AG mostró convergencia prematura: mejora rápido en las primeras generaciones y luego se estanca. Estrategias como mayor mutación, mecanismos de diversidad o búsqueda local podrían mejorar los resultados.
\item No es posible afirmar que el AG alcanzó el óptimo, porque el solver no lo demostró. El óptimo está entre 30\,817 y 48\,878.
\end{enumerate}

\subsection*{Limitaciones y trabajo futuro}
El estudio usa una sola instancia, un solo solver y un límite de 600 segundos. Trabajo futuro: dar más tiempo al solver o usar la solución del AG como punto de partida del método exacto, probar otros operadores de cruce y de selección, incorporar búsqueda local, y evaluar el método en más instancias, por ejemplo las de la OR-Library, con pruebas estadísticas entre configuraciones.
